%%%%%%%%%%%%
%
% $Autor: Wings $
% $Datum: 2019-03-05 08:03:15Z $
% $Pfad: SafetyGuidelines.tex $
% $Version: 4250 $
% !TeX spellcheck = en_GB/de_DE
% !TeX encoding = utf8
% !TeX root = manual 
% !TeX TXS-program:bibliography = txs:///biber
%
%%%%%%%%%%%%

\chapter{Safety Guidelines}

\section{Autonomous Operation Safety}

\begin{itemize}
    \item Always supervise autonomous runs; be prepared to stop the robot manually
    \item Establish safe boundaries and clear obstacles from test area
    \item Never combine autonomous mode with foot traffic in same area
    \item Use speed limiter on first runs in unfamiliar environments
\end{itemize}

\section{Testing Precautions}

\begin{itemize}
    \item Start with low motor speeds and gradually increase
    \item Test collision avoidance in controlled environment before deployment
    \item Do not rely solely on autonomous features; maintain manual override capability
    \item Report any unsafe behaviors immediately and disable autonomous mode
\end{itemize}

\section{Environmental Hazards}

\begin{itemize}
    \item Avoid stairs, ledges, and elevated surfaces unless tested
    \item Keep away from water sources
    \item Test on surfaces similar to intended deployment (carpet, tile, concrete)
    \item Be aware of differences between testing and real environments
\end{itemize}

\section{Emergency Procedures}

\begin{itemize}
    \item If robot malfunctions during testing, disconnect power immediately
    \item Document any near-miss incidents for analysis
    \item Do not attempt repairs while system is powered
\end{itemize}





